% =============================================================================
% Бүлэг 3: Өгөгдөл боловсруулалт
% =============================================================================

\chapter{Өгөгдөл боловсруулалт}
\label{ch:data}
\pagecolor{white}

Энэ бүлэгт SQL үүсгэх загвар ба орчуулгын загварт ашигласан өгөгдлийн эх сурвалж, боловсруулалт, чанарын хяналт, үнэлгээний
сангийн хуваалтыг тайлбарлана.

\section{SQL үүсгэх загварын өгөгдөл}
\label{sec:db}

\textbf{Схем.} Өмнөх ажлын репозиторт \cite{text2sql_thesis_2026} байсан ББСБ-ийн PostgreSQL схемийг хуулж ашиглав. 13
бизнес схемд нийт 189 хүснэгт ба view бий: байгууллага, харилцагч, KYC, бүтээгдэхүүн, зээл, барьцаа, данс ба гүйлгээ,
хадгаламж, барагдуулалт, эрсдэл, ерөнхий дэвтэр, тайлан, систем. Ачаалах үед илэрсэн алдааг зөвхөн хуулбар дээр засаж
бүртгэсэн; эх репозиторт өөрчлөлт хийгээгүй.

\textbf{Синтетик өгөгдөл.} Бодит харилцагчийн мэдээлэл ашиглах эрхгүй тул Python-оор детерминист (seed 42) синтетик
өгөгдөл үүсгэв. «Өнөөдөр» нь 2026-06-30, түүх 2024-01-01-ээс эхэлнэ: 3 000 харилцагч, 2 666 зээл, 212 566 гүйлгээ, 17
салбар, 158 ажилтан; нийт 648 205 мөр. Зээлийн симуляц нь хуваарь, төлөлт, хугацаа хэтрэлт, бүтцийн өөрчлөлт, шүүхийн
шат, хасалт, ерөнхий дэвтрийн бичилтийг хамарна.

\textbf{text2sql давхарга.} Загварт 189 объектыг шууд үзүүлэхийн оронд 28 view бүхий \tsql{t2s} схем бүтээв. View бүр нэг
бизнес объектод тохирно (\tsql{loans}, \tsql{loan\_payments}, \tsql{collection\_cases} г.м.). Код бүрийн хажууд нэрийн
багана (\tsql{status\_code} → \tsql{status\_name}), төгрөгт хөрвүүлсэн дүн (\tsql{*\_mnt}) нэмж, хувийн мэдээллийн
багануудыг (регистр, утас, хаяг) хассан.

\textbf{Мэдлэгийн сан.} Хүснэгт~\ref{tab:knowledge}. Схемийн тайлбар DDL-ийн comment хэлбэрээр prompt-д орно (Эх
код~\ref{lst:ddl}). Жишээний санг 160 → 355 → 635 болгон гурван шатаар өргөтгөв: LLM agent-уудын бичсэн жишээнээс
давхардал ба үнэлгээний асуулттай төстэйг хасаж, бүгдийг DB дээр гүйцэтгэж шалгасан (алдаа 0).

\begin{table}[H]
\centering\small
\caption{Мэдлэгийн сангийн бүрэлдэхүүн}
\label{tab:knowledge}
\begin{tabularx}{\textwidth}{p{3.4cm}p{3.0cm}X}
\toprule
Бүрэлдэхүүн & Хэмжээ & Агуулга \\
\midrule
Semantic model & 28 view, 425 багана & Монгол ба англи тайлбар, синоним, enum кодын нэр, гадаад түлхүүр \\
Нэр томьёо (glossary) & 45 (+52 нэмэлт) & Тодорхойлолт, каноник SQL, хамаарах хүснэгт; 7 нь хоёрдмол \\
KPI-ийн давхарга & 20 үзүүлэлт, 14 хэмжээс & PAR30, NPL хувь, барагдуулалтын хувь г.м. \\
Утгын индекс & багана бүрт ≤400 утга & Хувийн мэдээллээс бусад текст багана; галигийн «сул түлхүүр» \\
Жишээний сан & 160 → 355 → 635 & Асуулт ба SQL \\
\bottomrule
\end{tabularx}
\end{table}

\begin{lstlisting}[language=SQLmn, basicstyle=\ttfamily\scriptsize, caption={Тайлбартай DDL-ийн хэсэг (prompt-д орох хэлбэр)}, label={lst:ddl}]
CREATE TABLE loan_applications (
  application_id TEXT,  -- (*@Хүсэлтийн ID@*)
  status_code TEXT,  -- (*@Хүсэлтийн төлөв;@*) SUBMITTED=(*@Илгээгдсэн@*), REJECTED=(*@Татгалзсан@*), ...
  channel TEXT,  -- (*@Хүсэлт ирсэн суваг;@*) BRANCH=(*@Салбар@*), MOBILE=(*@Мобайл апп@*), WEB=(*@Вэб@*), ...
  requested_amount NUMERIC,  -- (*@Хүссэн дүн@*)
  ...
\end{lstlisting}

\textbf{Нэмэлт сургалтын өгөгдөл.} \textbf{C бүлэг}: BIRD train-ийн 5 000 асуултыг Qwen3.5-9B-ээр монгол руу орчуулж (SQL-ийн
утга, нэрийг маскаар хамгаалж), төсвийн улмаас 2 000 асуултаар rejection sampling хийв (асуулт бүрт 8 хариу); C2 нь
gold-оор чиглүүлсэн reasoning, C3 нь өөрийн DB-ийн синтетик асуулт ба 160 жишээг нэмсэн. \textbf{D2}: B5-ийн RAG prompt-тай
3 756 мөр; өөрийн DB-ийн 344 жишээ бүгд prompt-ийн урт 16 384 токеноос хэтэрсэн тул \textbf{хасагдсан}. \textbf{D3}: schema
linking-ээр prompt-ийн медиан уртыг 5 546 токен болгож (gold баганын recall 99.0\%), өөрийн DB-ийн жишээг сургалтад
багтаасан. \textbf{S1--S4}: схем, тодорхойлолт өөрчлөгдөхөд тогтвортой байх сургалтын prompt --- эх схем (629), нэр сольсон
схем (2 516), өөрчилсөн тодорхойлолт (64), татгалзах асуулт (480).

\section{Орчуулгын загварын өгөгдөл}
\label{sec:mt-data}

Алдааны шинжилгээгээр орчуулга гол алдааны эх үүсвэр болсон тул (§\ref{sec:errors}) санхүүгийн хэллэгийн орчуулгын
өгөгдөл бэлтгэв (Хүснэгт~\ref{tab:mt-data}). Гаднаас өгөгдөл татаагүй; бүгдийг LLM agent-ууд бичиж, өөр agent сохроор
хянасан (123 засвар). Латин галигийн хоёр хэв маягийг (ү/ө-г v/u эсвэл u/o) хамруулсан.

\begin{table}[H]
\centering\small
\caption{Орчуулгын өгөгдөл (MT v2)}
\label{tab:mt-data}
\begin{tabularx}{\textwidth}{p{5.0cm}X}
\toprule
Бүрэлдэхүүн & Хэмжээ \\
\midrule
Асуултын хос (монгол → англи) & 2 828 хос (val 287), 8 домэйн \\
Санхүүгийн өвөрмөц хэллэг & 123 хэллэг + жишээ асуулт \\
Минимал хос (нэг үгээр ялгаатай) & 256 хос: нэр томьёо, босго, хугацаа, тоо, эрэмбэ, үгүйсгэл \\
Бичгийн хувилбар & кирилл, латин (хэрэглэгчийн ба синтетик), холимог, англи → монгол \\
Сургалтын мөр (SFT) & 8 601 мөр (MADLAD-ийн формат: \tsql{<2en> [T: нэр томьёо] асуулт}) \\
Хүний шалгалтын хуудас & 150 хос (бэлэн, шалгагдаагүй) \\
\bottomrule
\end{tabularx}
\end{table}

\section{Өгөгдлийн нэгдсэн үзүүлэлт ба хуваалт}
\label{sec:evalset}

\textbf{Үнэлгээний сан.} Монгол Text2SQL-д бэлэн үнэлгээний сан байхгүй тул ББСБ-ийн 8 домэйнд асуулт бэлтгэв. Асуулт бүр
кирилл, латин галиг, англи (заримдаа холимог) хувилбар; ангилал (хариулах боломжтой, боломжгүй, хоёрдмол); хүндрэл; gold
ба орлох SQL; хоёрдмол асуултын тайлбарууд; хариулах боломжгүй шалтгаантай. Gold SQL бүрийг үндсэн DB ба өгөгдлийг
өөрчилсөн 8 test-suite хуулбар дээр гүйцэтгэж, бие даасан 4 хянагч agent утга, ангилал, давхардлыг шалгав.

\textbf{Хуваалт.} Ижил загвараас үүссэн асуултууд (template family) нэг талд байхаар dev ба test-д хуваав
(Хүснэгт~\ref{tab:evalset}); жишээ асуултыг Хүснэгт~\ref{tab:eval-examples}-д, чанарын шалгалтыг
Хүснэгт~\ref{tab:quality}-д үзүүлэв.

\begin{table}[H]\centering\small
\caption{Үнэлгээний сангийн хуваалт ба ангилал}\label{tab:evalset}
\begin{adjustbox}{max width=\textwidth}
\begin{tabular}{lrrrrrrr}\toprule
Хуваалт & Нийт & Хариулах & Хариулах & Хоёрдмол & Хүндрэл & Холимог & Загварын \\
 & & боломжтой & боломжгүй & & хялбар/дунд/хүнд & хувилбар & гэр бүл \\\midrule
dev & 300 & 224 & 38 & 38 & 57/121/46 & 56 & 271 \\
test & 200 & 151 & 24 & 25 & 38/82/31 & 38 & 177 \\
\bottomrule\end{tabular}
\end{adjustbox}
\end{table}

\begin{table}[H]
\centering\scriptsize
\caption{Үнэлгээний сангийн жишээ асуултууд}
\label{tab:eval-examples}
\begin{tabularx}{\textwidth}{>{\raggedright\arraybackslash}p{1.6cm}L>{\raggedright\arraybackslash}p{4.2cm}}
\toprule
Ангилал & Асуулт (кирилл / галиг) & Хүлээгдэх үр дүн \\
\midrule
Хялбар & 2025 онд нийт хэдэн зээлийн хүсэлт ирсэн бэ? / \emph{2025 ond niit heden zeeliin hvselt irsen be?} &
  \tsql{count(*) ... BETWEEN '2025-01-01' AND '2025-12-31'} \\
Хүнд & 2021 оноос хойш он бүрээр ажилд орсон болон ажлаас гарсан ажилтны тоог зэрэгцүүлэн гарга. &
  Хоёр CTE + он бүрээр нэгтгэл \\
Холимог & Энэ оны эхнээс mobile app-аар хэдэн loan application ирэв? & \tsql{channel = 'MOBILE'}, 2026-01-01 … 2026-06-30 \\
Боломжгүй & Хар жагсаалтад орсон харилцагчдын регистрийн дугаар, гэрийн хаягийг жагсааж өгөөч. &
  Татгалзах (хувийн мэдээлэл) \\
Хоёрдмол & Аль бүтээгдэхүүн дээр хамгийн их татгалзсан бэ? & Тодруулах: тоогоор уу, хувиар уу? \\
\bottomrule
\end{tabularx}
\end{table}

\begin{table}[H]
\centering\small
\caption{Өгөгдлийн чанарын нэгдсэн үзүүлэлт}
\label{tab:quality}
\begin{tabularx}{\textwidth}{X p{5.0cm}}
\toprule
Шалгалт & Үр дүн \\
\midrule
Gold SQL: 500 асуулт × (үндсэн DB + 8 test-suite хуулбар) & 0 алдаа; oracle EX 100\% \\
Dev--test ба жишээ--үнэлгээний ойролцоо давхардал & 0 \\
Мэдлэгийн сан ↔ DB (425 багана, 45 нэр томьёо) & 0 зөрчил \\
Жишээний сан (635): DB дээр гүйцэтгэх & 0 алдаа \\
MT v2: сохор хяналт & 123 засвар \\
Нэр сольсон схем (S2): статик шалгалт & 100\% \\
Хүний (домэйн шинжээчийн) шалгалт & \textbf{хийгдээгүй} (хуудас бэлэн) \\
\bottomrule
\end{tabularx}
\end{table}

\section{Бүлгийн дүгнэлт}
\label{sec:ch4-summary}

Хувийн мэдээлэлгүй синтетик өгөгдлийн сан, 28 view бүхий text2sql давхарга, монгол мэдлэгийн сан, загварын гэр бүлээр
хуваасан 500 асуултын үнэлгээний сан, SQL бичигч ба орчуулагчийн сургалтын өгөгдлийг бэлтгэв. Гол сул тал нь бүх асуулт,
жишээг LLM agent-ууд бичиж хянасан бөгөөд хүний шалгалт хийгдээгүй явдал юм.
